\chapter{Estudio de Seguridad y Salud}
\label{ch:sys}
\section{Objetivo}
\label{sec:sys-objetivo}
El objeto del presente Estudio de Seguridad y Salud es la redacción de los documentos necesarios que definan, en el marco del Real Decreto 1627/1997, de 24 de Octubre, por el que se establecen disposiciones mínimas de seguridad y salud en las obras de construcción, las previsiones y desarrollo de las soluciones necesarias para los problemas de ejecución de la obra, y la prevención de riesgos de accidentes preceptivas de sanidad, higiene y bienestar de los trabajadores durante el desarrollo de la misma.
	
En aplicación de este Estudio de Seguridad y Salud de la obra, cada contratista, subcontratista y trabajadores autónomos, elaborarán un plan de seguridad y salud en el trabajo, en el que se analicen , estudien, desarrollen y complementen las previsiones contenidas en este estudio.
\section{Datos generales de la obra}
\label{sec:sys-obra}
El presente Estudio Básico de Seguridad y Salud se refiere al Proyecto de la línea aérea de alta tensión, cuyos datos generales son:
\begin{itemize}[label=--]
    \item \textit{Proyecto de Ejecución:} Línea de A.T. L'Albir-Benidorm.
    \item \textit{Autor del Proyecto:} Mario Carrión Sirvent
    \item \textit{Emplazamiento:} Alicante
    \item \textit{Presupuesto de Ejecución material:} \SI{65670}{€}.
\end{itemize}
Las unidades constructivas que componen la presente obra son:
\begin{itemize}
    \item Replanteo.
    \item Desbroce.
    \item Excavación.
    \item Cimentación.
    \item Armado e izado de apoyos.
    \item Instalación de conductores desnudos.
    \item Instalación de aisladores
    \item Instalación de crucetas.
    \item Instalación de aparatos de seccionamiento y corte (interruptores, seccionadores, fusibles...).
    \item Instalación de limitadores de sobretensión (autoválvulas).
    \item Instalación de transformadores tipo intemperie sobre apoyos.
    \item Instalación de dispositivos antivibraciones.
    \item Medida de altura de conductores.
    \item Detección de partes en tensión.
    \item Interconexión entre elementos.
    \item Conexión y desconexión de líneas o equipos.
    \item Puesta a tierra y conexiones equipotenciales.
\end{itemize}
\section{Normativa aplicable}
\label{sec:sys-normativa}
\subsection{Normas oficiales}
\label{subsec:sys-normasOficiales}
Son de obligado cumplimiento todas las disposiciones legales o reglamentarias, resoluciones y cuantas otras fuentes normativas contengan concretas regulaciones en materia de Seguridad e Higiene en el trabajo, propias de la Industria Eléctrica o de carácter general, que se encuentren vigentes y sean de aplicación durante el tiempo en el que subsista la relación contractual promotor-contratista, según las actividades a realizar.

En particular:
\begin{itemize}
    \item Ley 8/1980, de 1 de marzo, del Estatuto de los Trabajadores
    \item Ordenanza General de Seguridad e Higiene en el Trabajo (9 de marzo de 1971).
    \item Homologación de medios de Protección personal de los trabajadores (BOL. de 29 de mayo de 1974. Orden de 15 de julio de 1974).
    \item Estatuto de los Trabajadores (Ley 8/1980, de 20 de marzo).
    \item Ley de Prevención de Riesgos Laborales (Ley 31/1995, de 8 de noviembre).
    \item Real Decreto 39/1997, de 17 de enero, por el que se aprueba el Reglamento de los Servicios de Prevención.
    \item Orden de 27 de junio de 1997, por la que se desarrolla el RD 39/1.997, de 17 de enero.
    \item Real Decreto 485/1997, de 14 de abril, sobre disposiciones mininas en materia de señalización de seguridad y salud en el trabajo.
    \item Real Decreto 486/1997, de 14 de abril, por el que se establecen disposiciones mínimas de seguridad y salud en los lugares de trabajo.
    \item Real Decreto 487/1997, de 14 de abril, sobre disposiciones mínimas de seguridad y salud relativas a la manipulación manual de cargas que entrañen riesgos, en particular dorso-lumbares, para los trabajadores.
    \item  Real Decreto 773/1997, de 30 de mayo, sobre disposiciones mínimas de seguridad y salud relativas a la utilización por los trabajadores de equipos de protección individual.
    \item Real Decreto 949/1997, de 20 de Junio, por el que se establece el certificado de profesionalidad de la ocupación de prevencionista de riesgos laborales. 
    \item Real Decreto 1215/1997, de 18 de julio, por el que se establecen las disposiciones mínimas de seguridad y salud para la utilización por los trabajadores de los equipos de trabajo.
    \item Real Decreto 1627/1997, de 24 de octubre, por el que se establecen disposiciones mínimas de seguridad y de salud en las obras de construcción.
    \item Real Decreto 223/2008 de 15 de febrero, por el que se aprueban el Reglamento sobre condiciones técnicas y garantías de seguridad en líneas eléctricas de alta tensión y sus instrucciones técnicas complementarias ITC-LAT 01 a 09.
    \item Reglamento sobre Condiciones Técnicas y de Garantía de Seguridad en Centrales Eléctricas, Subestaciones y Centros de transformación (Decreto 3275/1982 de 12 de noviembre) e instrucciones Técnicas Complementarias.
\end{itemize}
\subsection{Normas específicas}
\label{subsec:sys-normasEspecificas}
Dentro de estas Normas deben tener especialmente en cuenta todas las Recomendaciones, Prescripciones e Instrucciones de la Asociación de Medicina y Seguridad en el Trabajo de UNESA para la Industria Eléctrica (AMYS), que se recogen en:
\begin{itemize}
\item “Prescripciones de Seguridad para trabajos y maniobras en instalaciones eléctricas”.
\item “Prescripciones de Seguridad para trabajos mecánicos y diversos”.
\item Instrucción General para la realización de los trabajos en tensión en Alta Tensión y sus Desarrollos.
\item Instrucción General para la realización de los trabajos en tensión en Baja Tensión y sus Desarrollos.
\end{itemize}
\section{Obligación del promotor}
\label{sec:sys-promotor}
El promotor está obligado a incluir el presente Estudio de Seguridad y Salud, como documento del Proyecto de Obra.

Antes del inicio de los trabajos, designará un coordinador en materia de seguridad y salud, cuando en la ejecución de las obras intervengan más de una empresa, o empresas y trabajadores autónomos, o diversos trabajadores autónomos.

La designación de coordinadores en materia de seguridad y salud no eximirá al promotor de sus responsabilidades.

El promotor deberá efectuar un aviso a la autoridad laboral competente antes del comienzo de las obras, que se redactará con arreglo a lo dispuesto en el Anexo III del R.D. 1627/1997, de 24 de octubre, debiendo exponerse en la obra de forma visible y actualizándose si fuera necesario.
\section{El coordinador}
\label{sec:sys-coordinador}
El Coordinador en materia de seguridad y salud durante la ejecución de la obra, deberá coordinar los principios generales de prevención y de seguridad, tomando las decisiones técnicas y de organización con el fin de planificar los distintos trabajos o fases que vayan a desarrollarse simultánea o sucesivamente.
	
Deberá coordinar las actividades de la obra para garantizar que los contratistas y, en su caso, los subcontratistas y los trabajadores autónomos, apliquen de manera coherente y responsable los principios de la acción preventiva que se recogen en el artículo 15 de la Ley de prevención de Riesgos Laborales durante la ejecución de la obra y, en particular, en las tareas o actividades a que se refiere el artículo 10 del Decreto 1627/1997 de 24 de octubre, sobre disposiciones mínimas de seguridad y de salud en las obras de construcción.
	
El Coordinador deberá aprobar el Plan de Seguridad y Salud elaborado por el contratista y, en su caso, las modificaciones introducidas en el mismo.

Así mismo organizará la coordinación de actividades empresariales previstas en el artículo 24 de la Ley de Prevención de Riesgos Laborales y coordinará las acciones y funciones de control de la aplicación correcta de los métodos de trabajo.

El Coordinador deberá adoptar las medidas necesarias para que sólo las personas autorizadas puedan acceder a la obra.
\section{Contratistas y subcontratistas}
\label{sec:sys-contratistas}
Estarán obligados a aplicar los principios de la acción preventiva que se recogen en el artículo 15 de la Ley de Prevención de Riesgos Laborales, cumplir y hacer cumplir a su personal lo establecido en el Plan de Seguridad y Salud e informar y proporcionar las instrucciones adecuadas a los trabajadores autónomos sobre todas las medidas que hayan de adoptarse en lo que se refiere a seguridad y salud en la obra.

Deberán atender las indicaciones y cumplir las instrucciones del coordinador en materia de seguridad y salud durante la ejecución de la obra.

Los contratistas y subcontratistas serán responsables de la ejecución correcta de las medidas preventivas fijadas en el plan de seguridad y salud en lo relativo a las obligaciones que les correspondan a ellos directamente o, en su caso, a los trabajadores autónomos por ellos contratados.

Además los contratistas y subcontratistas responderán solidariamente de las consecuencias que se deriven del incumplimiento de las medidas previstas en el plan en los términos del apartado 2 del artículo 42 de la Ley de Prevención de Riesgos Laborales.

Las responsabilidades de los coordinadores, de la dirección facultativa y del promotor no eximirán de sus responsabilidades a los contratistas y a los subcontratistas.

Los equipos de protección individual a disponer para cada uno de los puestos de trabajo a desempeñar, determinadas en el Plan de Seguridad y Salud en el Trabajo a elaborar por el contratista, estarán en consonancia con el resultado previsto por éste en la evaluación de los riesgos que está obligado a realizar en cumplimiento del R.D. 39/1.997, de 17 de Enero, por el que se aprueba el Reglamento de los Servicios de Prevención. Una copia de dicha evaluación y de su resultado, se adjuntará al Plan en el momento de su presentación.

Asimismo, y en aplicación del R.D. 773/1.997, de 30 de Mayo, sobre disposiciones mínimas de seguridad y salud relativas a la utilización por los trabajadores de los equipos de protección individual, es responsabilidad del contratista suministrar dichas protecciones individuales a los trabajadores de manera gratuita, reponiéndolas cuando resulte necesario, motivo por el cual, dentro del Plan de Seguridad y Salud en el Trabajo a elaborar por el contratista, éstas se relacionarán exhaustivamente en todos los apartados del mismo, de acuerdo con lo señalado en el párrafo anterior, pero no se valorarán dentro del presupuesto del plan.
\section{Obligaciones de los trabajadores}
\label{sec:sys-obligaciones}
Los trabajadores están obligados a:
\begin{enumerate}
    \item Aplicar los principios de la acción preventiva que se recoge en el artículo 15 de la Ley de Prevención de Riesgos Laborales, y en particular:
    \begin{itemize}[label=--]
        \item Mantenimiento de la obra en buen estado de orden y limpieza.
        \item Almacenamiento y evacuación de residuos y escombros.
        \item Recogida de materiales peligrosos utilizados.
        \item Adaptación del periodo de tiempo efectivo que habrá de dedicarse a los distintos trabajos o fases de trabajo.
        \item Cooperación entre todos los intervinientes en la obra.
        \item Interacciones o incompatibilidades con cualquier otro trabajo o actividad.
    \end{itemize}
  \item Cumplir las disposiciones mínimas establecidas en el Anexo IV del R.D. 1627/1997.
  \item Ajustar su actuación conforme a los deberes sobre coordinación de las actividades empresariales previstas en le artículo 24 de la Ley de Prevención de Riesgos Laborales, participando en particular en cualquier medida de actuación coordinada que se hubiera establecido.
  \item Cumplir con las obligaciones establecidas para los trabajadores en el artículo 29, apartados 1 y 2 de la Ley de Prevención de Riesgos Laborales.
  \item Utilizar equipos de trabajo que se ajusten a lo dispuesto en el R.D. 1215/1997.
  \item Elegir y utilizar equipos de protección individual en los términos previstos en el R.D. 773/1997.
  \item Atender las indicaciones y cumplir las instrucciones del coordinador en materia de seguridad y salud.
\end{enumerate}
Los trabajadores autónomos deberán cumplir lo establecido en el plan de seguridad y salud.
\section{Libro de incidencias}
\label{sec:sys-incidencias}
En cada centro de trabajo existirá, con fines de control y seguimiento del plan de seguridad y salud, un libro de incidencias que constará de hojas duplicadas y que será facilitado por el colegio profesional al que pertenezca el técnico que haya aprobado el plan de seguridad y salud.

Deberá mantenerse siempre en obra y en poder del coordinador. Tendrán acceso al libro, la Dirección Facultativa, los contratistas y subcontratistas, los trabajadores autónomos, las personas con responsabilidades en materia de prevención de las empresas intervinientes, los representantes de los trabajadores, y los técnicos especializados de las Administraciones Públicas competentes en esta materia, quienes podrán hacer anotaciones en el mismo.

Efectuada una anotación en el libro de incidencias, el coordinador estará obligado a remitir en el plazo de \SI{24}{\hour} una copia a la Inspección de Trabajo y Seguridad Social de la provincia en que se realiza la obra. Igualmente notificará dichas anotaciones al contratista y a los representantes de los trabajadores.
\section{Derecho de los trabajadores}
\label{sec:sys-derechos}
Los contratistas y subcontratistas deberán garantizar que los trabajadores reciban una información adecuada y comprensible de todas las medidas que hayan de adoptarse en lo que se refiere a seguridad y salud en la obra.

Una copia del plan de seguridad y salud y de sus posibles modificaciones, a los efectos de su conocimiento y seguimiento, será facilitada por el contratista a los representantes de los trabajadores en el centro de trabajo.
\section{Prevención de riesgos profesionales}
\label{sec:sys-prl}
\subsection{Protecciones individuales generales}
\label{subsec:sys-prlEPI}
\begin{enumerate}
    \item Cascos: para todas las personas que participan en obra, incluidos visitantes.
    \item Guantes de uso general.
    \item Guantes de goma.
    \item Guantes de soldador.
    \item Guantes diacetílicos.
    \item Botas de agua.
    \item Botas de seguridad de lona.
    \item Botas de seguridad de cuero.
    \item Botas dialécticas.
    \item Gafas de soldador.
    \item Gafas de seguridad antiproyecciones.
    \item Pantalla de soldador.
    \item Mascarillas antipolvo.
    \item Protectores auditivos.
    \item Polainas de soldador.
    \item Manguitos de soldador.
    \item Mandiles de soldador.
    \item Cinturón de seguridad de sujeción.
    \item Cinturón antivibratorio.
    \item Chalecos reflectantes.
\end{enumerate}
\subsection{Protecciones colectivas generales}
\label{subsec:sys-prlEPC}

\begin{enumerate}
    \item Pórticos protectores de líneas eléctricas.
    \item Vallas de limitación y protección.
    \item Señales de seguridad.
    \item Cintas de balizamiento.
    \item Redes.
    \item Soportes y anclajes de redes.
    \item Tubo sujeción cinturón de seguridad.
    \item Anclaje para tubo.
    \item Balizamiento luminoso.
    \item Extintores.
    \item Interruptores diferenciales
    \item Toma de tierra.
    \item Válvula antiretroceso.
    \item Riegos.
\end{enumerate}
\subsection{Formación}
\label{subsec:sys-prlFormacion}
Todo personal debe recibir, al ingresar en la obra, una exposición de los métodos de trabajo y los riesgos que éstos pudieran entrañar, juntamente con las medidas de seguridad que deberá emplear.

Eligiendo al personal más cualificado impartirán cursillos de socorrismo y primeros auxilios, de forma que todos los trabajos dispongan de algún socorrista.

Se informará a todo el personal interviniente en la obra, sobre la existencia de productos inflamables, tóxicos, etc. y medidas a tomar en cada caso.
\subsection{Medicina preventiva y primeros auxilios}
\label{subsec:sys-prlMedicina}
Se tendrán en cuenta las siguientes consideraciones:
\begin{enumerate}
    \item Botiquín: Deberá existir en la obra al menos un botiquín con todos los elementos suficientes para curas, primeros auxilios, dolores, etc.
    \item Asistencia a accidentados: Se deberá informar a la obra del emplazamiento de los diferentes Centros Médicos, Residencia Sanitaria, médicos, ATS., etc., donde deba trasladarse a los posibles accidentados para un más rápido y efectivo tratamiento, disponiendo en la obra de las direcciones, teléfonos, etc., en sitios visibles.
    \item Reconocimiento Médico: todo el personal que empiece a trabajar en la obra deberá pasar un reconocimiento médico previo que certifique su aptitud.
    \item Instalaciones: se dotará a la obra, si así se estima en el correspondiente Plan de Seguridad, de todas las instalaciones necesarias, tales como:
    \begin{itemize}[label=--]
        \item Almacenes y talleres:
        \item Vesturarios y Servicios:
        \item Comedor o, en su defecto, locales particulares para el mismo fin:
    \end{itemize}
\end{enumerate}
\section{Identificación de riesgos y medidas preventivas a aplicar}
\label{sec:sys-riesgos}
\subsection{Fase de actuaciones previas}
\label{subsec:sys-riesgosActuaciones}
En esta fase se consideran las labores previas al inicio de las obras, como puede ser el replanteo, red de saneamiento provisional para vestuarios y aseos de personal de obra...
\subsubsection{Riesgos detectables}
\label{subsubsec:sysActuaciones-RD}
\begin{itemize}
    \item Atropellos y colisiones originados por maquinaria.
    \item Vuelcos y deslizamientos de vehículos de obra.
    \item Caídas en el mismo nivel.
    \item Torceduras de pies.
    \item Generación de polvo.
\end{itemize}
\subsubsection{Medidas de seguridad}
\label{subsubsec:sysActuaciones-MS}
\begin{itemize}
    \item Se cumplirá la prohibición de presencia de personal, en las proximidades y ámbito de giro de maniobra de vehículos y en operaciones de carga y descarga de materiales.
    \item La entrada y salida de camiones de la obra a la vía pública, será debidamente avisada por persona distinta al conductor.
    \item Será llevado un perfecto mantenimiento de maquinaría y vehículos.
    \item La carga de materiales sobre camión será correcta y equilibrada y jamás superará la carga máxima autorizada.
    \item El personal irá provisto de calzado adecuado.
    \item Todos los recipientes que contengan productos tóxicos o inflamables, estarán herméticamente cerrados.
    \item No se apilarán materiales en zonas de paso o de tránsito, retirando aquellos que puedan impedir el paso.
\end{itemize}
\subsubsection{Prendas de protección de personal}
\label{subsubsec:sysActuaciones-EPI}
\begin{itemize}
    \item Casco homologado.
    \item Mono de trabajo y en su caso, trajes de agua y botas de goma de media caña.
    \item Empleo de cinturones de seguridad por parte del conductor de la maquinaría si no está dotada de cabina y protección antivuelco.
    \item Mascarillas antipolvo con filtro mecánico.

\end{itemize}
\subsection{Fase de acopio de material}
\label{subsec:sys-riesgosMaterial}
\subsubsection{Riesgos detectables}
\label{subsubsec:sysMaterial-RD}
\begin{itemize}
    \item Caídas de objetos
    \item Golpes.
    \item Heridas
    \item Sobreesfuerzos.
\end{itemize}
\subsubsection{Medidas de seguridad}
\label{subsubsec:sysMaterial-MS}
\begin{itemize}
    \item Antes de comenzar el acopio de material a los lugares de trabajo, se deberá realizar un reconocimiento del terreno, con el fin de escoger la mejor ruta.
    \item En el caso en que para acceder al lugar de trabajo fuera necesario adecuar o construir una ruta de acceso, esta deberá realizarse con la maquinaria y medios adecuados.
\end{itemize}
\subsubsection{Prendas de protección de personal}
\label{subsubsec:sysMaterial-EPI}
\begin{itemize}
    \item Guantes comunes de trabajo de lona y piel flor.
    \item Ropa de trabajo cubriendo la mayor parte del cuerpo.
    \item Botas reforzadas.
\end{itemize}
\subsection{Carga y descarga de materiales}
\label{subsec:sys-riesgosCargaMat}
\subsubsection{Riesgos detectables}
\label{subsubsec:sysCarga-RD}
\begin{itemize}
    \item Caída de operarios al mismo nivel.
    \item Golpes, heridas y sobreesfuerzos.
    \item Caída de objetos.
\end{itemize}
\subsubsection{Medidas de seguridad}
\label{subsubsec:sysCarga-MS}
\begin{itemize}
    \item Con el fin de evitar posibles lesiones en la columna vertebral, el operario llevará a cabo el levantamiento de la carga realizando el esfuerzo con las piernas, y manteniendo en todo momento la columna recta.
    \item Un operario no podrá levantar más de 50 Kg en la carga y descarga manual. En el caso en concreto en que la carga fuera superior a la cantidad límite, se deberá realizar entre más trabajadores.
    \item En el caso en que el acarreo de pesos se estime en una duración superior a las 4 horas de trabajo continuadas, el peso máximo a acarrear será de 25 Kg., o bien deberán utilizarse medios mecánicos adecuados.
    \item Para la carga y descarga con medios mecánicos, la maquinaria a emplear deberá ser la adecuada (grúa, pala cargadora, etc.) y su maniobra deberá ser dirigida por personal especializado, no debiéndose superar en ningún momento la carga máxima autorizada.
    \item Todas las máquinas que participen en las operaciones deberán estar correctamente estabilizadas. La elevación de la carga deberá realizarse de forma suave y continuada.
    \item En el transcurso de operaciones de carga y descarga, ninguna persona ajena se acercará al vehículo. Debe acotarse el entorno y prohibirse el permanecer o trabajar dentro del radio de acción del brazo de una máquina
    \item Nunca permanecerá ni circulará personal debajo de las cargas suspendidas, ni permanecerá sobre las cargas.
    \item Para la descarga de bobinas de conductores, se emplearán cuerdas, rampas, raíles...
    \item Bajo ningún concepto se hará rodar la bobina por un solo canto.
    \item Se prohíbe el acopio de materiales a menos de 2 metros de las coronaciones de taludes.
\end{itemize}
\subsubsection{Prendas de protección de personal}
\label{subsubsec:sysCarga-EPI}
\begin{itemize}
    \item Guantes adecuados
    \item Ropa de trabajo.
    \item Botas de seguridad.
    \item Fajas antilumbago, si existen cargas muy pesadas.
\end{itemize}
\subsection{Movimientos de tierra y excavación}
\label{subsec:sys-riesgosExcavacion}
\subsubsection{Riesgos detectables}
\label{subsubsec:sysExcavacion-RD}
\begin{itemize}
    \item Choque, atropellos y atrapamientos ocasionados por la maquinaria.
    \item Vuelcos y deslizamientos de las máquinas.
    \item Caídas en altura del personal que intervienen en el trabajo.
    \item Generación de polvo.
    \item Desprendimiento de tierra y proyección de rocas.
    \item Caídas de personal al interior de pozos.
    \item Caídas a distinto nivel.
\end{itemize}
\subsubsection{Medidas de seguridad}
\label{subsubsec:sysExcavacion-MS}
\begin{itemize}
    \item En el caso de uso de herramientas, debido a las reducidas dimensiones que generalmente tendrán los hoyos, se recomienda que sea un único trabajador el que permanezca en su interior, para evitar accidentes por alcance entre ellos de las herramientas a emplear.
    \item Los picos, palas y otras herramientas deberán estar en buenas condiciones.
    \item En el caso de hoyos con probable peligro de derrumbamiento de paredes, nunca deberá quedar un operario solo en su interior, sino que en el exterior de hoyo debe permanecer, al menos, otro operario, para caso de auxilio.
    \item Las maniobras de las máquinas estarán dirigidas por persona distinta al conductor.
    \item Los escombros procedentes de la excavación deberán situarse a una distancia adecuada del hoyo, para evitar la caída al interior del mismo.
    \item Los pozos de cimentación se señalizarán para evitar caídas del personal a su interior desde su realización hasta que sean rellenados.
    \item Durante la ausencia de los operarios de la obra, los hoyos serán tapados con tablones u otros elementos adecuados.
    \item Se cumplirá la prohibición de presencia del personal en la proximidad de las máquinas durante su trabajo. 
    \item Durante la retirada de árboles no habrá personal trabajando en planos inclinados con fuerte pendiente.
    \item Mantenimiento correcto de la maquinaria.
    \item Al proceder a la realización de excavaciones, correcto apoyo de las máquinas excavadoras en el terreno.
    \item Si se realizan excavaciones de hoyos en roca que exijan uso de explosivos, la manipulación de estos deberá ser realizada por personal especializado, con el correspondiente permiso oficial y poseedor del carné de dinamitero.
    \item En caso de que sobrase dinamita, se entregará en el Cuartel de la Guardia Civil o se destruirá en obra.
\end{itemize}
\subsubsection{Prendas de protección de personal}
\label{subsubsec:sysExcavacion-EPI}
\begin{itemize}
    \item El equipo de los operarios que efectúen las labores de excavación estará formado por: ropa adecuada de trabajo, guantes adecuados, casco de seguridad, botas reforzadas y gafas antipolvo reforzadas si existiese la posibilidad de que pueda penetrar tierra y otras partículas en los ojos.
    \item Empleo del cinturón de seguridad por parte del conductor de la maquinaria.
\end{itemize}
\subsection{Cimentación}
\label{subsec:sys-riesgosCimentacion}
\subsubsection{Riesgos detectables}
\label{subsubsec:sysCimentacion-RD}
\begin{itemize}
    \item Caída de persona y/o objetos al mismo nivel.
    \item Caída de persona y/o objetos a distinto nivel.
    \item Contactos con el hormigón por salpicaduras en cara y ojos.
    \item Quemadura de la piel por la acción del cemento.
    \item Caída de la hormigonera por efecto del volteo por no estar suficientemente nivelada y sujeta.
\end{itemize}
\subsubsection{Medidas de seguridad}
\label{subsubsec:sysCimentacion-MS}
\begin{enumerate}[label=\alph*)]
    \item Vertidos directos mediante canaleta:
    \begin{itemize}
        \item Se instalarán fuertes topes de recorrido de los camiones hormigonera, para evitar vuelcos.
        \item Se prohíbe acerar las ruedas de los camiones hormigoneras a menos de 2 metros del borde de la excavación.
        \item Se prohíbe situar a los operarios detrás de los camiones hormigonera durante el retroceso.
        \item La maniobra de vertidos será dirigida por u capataz que vigilará que no se realicen maniobras inseguras.
    \end{itemize}
    \item Vertidos directos mediante cubo o cangilón:
    \begin{itemize}
        \item Se prohíbe cargar el cubo por encima de la carga máxima admisible de la grúa que lo sustenta.
        \item Se señalizará, mediante una traza horizontal ejecutada con pintura en color amarilla, el nivel máximo de llenado del cubo para no sobrepasar la carga admisible.
        \item La apertura del cubo para vertido se ejecutará exclusivamente accionando la palanca para ello, con las manos protegidas con guantes impermeables
        \item La maniobra de aproximación, se dirigirá mediante señales preestablecidas fácilmente inteligibles por el gruísta.
    \end{itemize}
\end{enumerate}
En general, habrá que tomar las siguientes medidas preventivas:
\begin{itemize}
    \item Ningún trabajador con antecedentes de problemas cutáneos participará en las labores de hormigonado.
    \item Si por alguna causa, algún trabajador sufriese lesiones por acción del cemento, se deberá notificar la aparición de las mismas lo antes posible, con el fin de evitar la cronificación y nuevas sensibilizaciones.
    \item Si el amasado se realiza con hormigonera in situ, ésta deberá estar correctamente nivelada y sujeta.
    \item Los trabajadores deberán tener especial cuidado con:
    \begin{itemize}
        \item No utilizar prendas con elementos colgantes y que no sean de la talla adecuada.
        \item No exponer la piel al contacto con el cemento.
        \item Realizar las operaciones con las debidas condiciones de estabilidad.
        \item No manejar elementos metálicos sin usar guantes adecuados.
        \item Utilizar el casco protector y gafas de protección si existe riesgo de que penetren partículas en los ojos.
    \end{itemize}
\end{itemize}
\subsubsection{Prendas de protección de personal}
\label{subsubsec:sysCimentacion-EPI}
\begin{itemize}
    \item Casco de seguridad.
    \item Gafas protectoras.
    \item Ropas y guantes adecuados.
    \item Faja antilumbago.
\end{itemize}
\subsection{Izado y armado de apoyos}
\label{subsec:sys-riesgosApoyos}
\subsubsection{Riesgos detectables}
\label{subsubsec:sysApoyos-RD}
\begin{itemize}
    \item Caída de personal desde altura.
    \item Atrapamientos.
    \item Golpes y heridas.
\end{itemize}
\subsubsection{Medidas de seguridad}
\label{subsubsec:sysApoyos-MS}
\begin{itemize}
    \item · No participarán en el armado de apoyos ningún operario con antecedentes de vértigo o epilepsia.
    \item Los desplazamientos de operarios por los apoyos se realizarán con las manos libres y siempre bien sujetos por el cinturón de seguridad.
    \item Se utilizarán grúas adecuadas (camión grúa, pluma...) según el peso y la altura, para el izado del apoyo. Cuidándose mucho de no sobrepasar la carga máxima autorizada.
    \item El manejo de la misma lo realizará siempre personal especializado.
    \item La grúa deberá estar en todo momento perfectamente nivelada.
    \item La elevación de las cargas deberá realizarse lentamente, evitando todo arranque o paro bruscos.
    \item Las maniobras deberán ser dirigidas por personal especializado, debiendo ser una única persona la encargada de dirigir al operador.
    \item En ningún momento deberá permanecer ninguna persona sobre las cargas ni sobre la maquinaria.
    \item La permanencia o circulación bajo carga suspendida queda terminantemente prohibida.
    \item Se tomarán especiales cuidados en la vestimenta cuando se trabaje con soldaduras.
    \item Una vez izado el apoyo deberá dejarse debidamente aplomado y estable.
    \item El armado del apoyo se realizará cuando el cimiento esté consolidado.
    \item Los apoyos sin hormigonar nunca se dejarán izados en ausencia de personal.
    \item Las herramientas y materiales no se lanzarán bajo ningún concepto, siempre se subirán y bajarán con la ayuda de cuerdas.
    \item Los trabajadores que realicen estos trabajos deberán usar cinturones portaherramientas.
\end{itemize}
\subsubsection{Prendas de protección de personal}
\label{subsubsec:sysApoyos-EPI}
\begin{itemize}
    \item Cascos de seguridad.
    \item Cinturón de seguridad que se amarrará a partes fijas de la torre.
    \item Ropas y guantes adecuados.
    \item Botas de seguridad.
\end{itemize}
\subsection{Montaje y apriete de tornillería}
\label{subsec:sys-riesgosTornilleria}
\subsubsection{Riesgos detectables}
\label{subsubsec:sysTornilleria-RD}
\begin{itemize}
    \item Cascos de seguridad.
    \item Cinturón de seguridad que se amarrará a partes fijas de la torre.
    \item Ropas y guantes adecuados.
    \item Botas de seguridad.
\end{itemize}
\subsubsection{Medidas de seguridad}
\label{subsubsec:sysTornilleria-MS}
\begin{itemize}
    \item Se utilizarán herramientas adecuadas, según el esfuerzo que haya que realizar, para el apriete de los tornillos.
    \item En el trabajo de apriete de tornillería trabajarán como máximo dos operarios, situados al mismo nivel o a trebolillos, y siempre en la cara externa del apoyo.
    \item La subida y bajada de material y herramientas se realizará con la ayuda de cuerdas, nunca lanzándolas.
    \item Los desplazamientos de los operarios por el apoyo se realizará con las manos libres y cinturón de seguridad.
\end{itemize}
\subsubsection{Prendas de protección de personal}
\label{subsubsec:sysTornilleria-EPI}
\begin{itemize}
    \item Cascos de seguridad.
    \item Cinturón de seguridad que se amarrará a partes fijas de la torre.
    \item Ropas y guantes adecuados.
    \item Botas de seguridad.
\end{itemize}
\subsection{Colocación de herrajes y aisladores. Tendido, tensado y engrapado de conductores}
\label{subsec:sys-riesgosHerrajes}
\subsubsection{Riesgos detectables}
\label{subsubsec:sysHerrajes-RD}
\begin{itemize}
    \item Caída de personal desde altura.
    \item Caídas de objetos desde altura.
    \item Golpes y heridas.
\end{itemize}
\subsubsection{Medidas de seguridad}
\label{subsubsec:sysHerrajes-MS}
\begin{itemize}
    \item Estas labores serán realizadas por personal especializado.
    \item El personal realizará su trabajo siempre con cinturón de seguridad sujeto a las partes fijas del apoyo y con la manos libres.
    \item Se entenderán la zona interior de los apoyos y las proyecciones de las crucetas como zonas peligrosas.
    \item Los gatos que soporten las bobinas dispondrán de elementos de frenado que impidan el movimiento rotatorio de la bobina.
    \item Las poleas de tendido deberán amarrarse adecuadamente a las cadenas de aisladores.
    \item En las operaciones de tensado y flechado, los apoyos fin de línea deberán estar arriostrados, de manera que no sufran esfuerzos superiores a los previstos en las condiciones normales de trabajo.
    \item Durante las operaciones de tendido y tensado el operario no deberá permanecer dentro del radio de acción del conductor.
    \item Para efectuar correctamente estas operaciones se usarán aparatos radioteléfonos, y de esta manera transmitir todas las órdenes de parada y puesta en marcha del tendido, o poner el alerta de cualquier imprevisto.
    \item Con el fin de evitar las descompensación de las crucetas, el flechado se realizará alternativamente en cada cruceta.
    \item Si fuera necesario, en los cruces con carreteras, ríos, calles, otras líneas... se instalarán protecciones (pórticos), según el tipo de cruzamiento, con el fin de proteger la zona de cruce, con el fin de evitar daños a terceros.
    \item Los cables se procurará pasarlos sobre cualquier obstáculo existente, de esta manera se evitarán resistencias a la hora de realizar el tendido.
\end{itemize}
\subsubsection{Prendas de protección de personal}
\label{subsubsec:sysHerrajes-EPI}
\begin{itemize}
    \item Cascos de seguridad.
    \item Cinturón de seguridad.
    \item Ropas y guantes adecuados.
    \item Botas de seguridad.
    \item Cinturón antilumbago.
\end{itemize}
\subsection{Uso de maquinarias y herramientas}
\label{subsec:sys-riesgosMaquinaria}
\subsubsection{Riesgos detectables}
\label{subsubsec:sysMaquinaria-RD}
\begin{itemize}
    \item Caída de personal desde altura.
    \item Caídas de objetos desde altura.
    \item Golpes y heridas.
\end{itemize}
\subsubsection{Medidas de seguridad}
\label{subsubsec:sysMaquinaria-MS}
\begin{itemize}
    \item Estas labores serán realizadas por personal especializado.
    \item El personal realizará su trabajo siempre con cinturón de seguridad sujeto a las partes fijas del apoyo y con la manos libres.
    \item Se entenderán la zona interior de los apoyos y las proyecciones de las crucetas como zonas peligrosas.
    \item Los gatos que soporten las bobinas dispondrán de elementos de frenado que impidan el movimiento rotatorio de la bobina.
    \item Las poleas de tendido deberán amarrarse adecuadamente a las cadenas de aisladores.
    \item En las operaciones de tensado y flechado, los apoyos fin de línea deberán estar arriostrados, de manera que no sufran esfuerzos superiores a los previstos en las condiciones normales de trabajo.
    \item Durante las operaciones de tendido y tensado el operario no deberá permanecer dentro del radio de acción del conductor.
    \item Para efectuar correctamente estas operaciones se usarán aparatos radioteléfonos, y de esta manera transmitir todas las órdenes de parada y puesta en marcha del tendido, o poner el alerta de cualquier imprevisto.
    \item Con el fin de evitar la descompensación de las crucetas, el flechado se realizará alternativamente en cada cruceta.
    \item Si fuera necesario, en los cruces con carreteras, ríos, calles, otras líneas... se instalarán protecciones (pórticos), según el tipo de cruzamiento, con el fin de proteger la zona de cruce, con el fin de evitar daños a terceros.
    \item Los cables se procurará pasarlos sobre cualquier obstáculo existente, de esta manera se evitarán resistencias a la hora de realizar el tendido.
\end{itemize}
\subsubsection{Prendas de protección de personal}
\label{subsubsec:sysMaquinaria-EPI}
\begin{itemize}
    \item Cascos de seguridad.
    \item Cinturón de seguridad.
    \item Ropas y guantes adecuados.
    \item Botas de seguridad.
    \item Cinturón antilumbago.
    \item Protección auditiva en caso necesario.
\end{itemize}
\section{Instalación eléctrica provisional en obra}
\label{sec:sys-instalacion}
El montaje de aparatos eléctricos será ejecutado por personal especialista, en prevención de los riesgos por montajes incorrectos.
	
El calibre o sección del cableado será siempre el adecuado para la carga eléctrica que ha de soportar.

Los hilos tendrán la funda protectora aislante sin defectos apreciables (rasgones, repelones y asimilables). No se admiten tramos defectuosos.

La distribución general, desde el cuadro general de la obra a los cuadros secundarios, se efectuará mediante manguera eléctrica antihumedad.

El tendido de los cables y mangueras, se efectuará a una altura mínima de 2 m. en los lugares peatonales y de 5 m. en los de vehículos, medidos sobre el nivel del pavimento.

Los empalmes provisionales entre mangueras, se ejecutarán mediante conexiones normalizadas estancas antihumedad.

Los interruptores se instalarán en el interior de cajas normalizadas, provistas de puerta de entrada con cerradura de seguridad.

Los cuadros eléctricos metálicos tendrán la carcasa conectada a tierra.

Los cuadros eléctricos se colgarán pendientes de tableros de madera recibidos a los paramentos verticales o bien a “pies derechos “firmes.

Las maniobras a ejecutar en el cuadro eléctrico general se efectuarán subido a una banqueta de maniobra o alfombrilla aislante.

Los cuadros eléctricos poseerán tomas de corriente para conexiones normalizadas blindadas para intemperie.

La tensión siempre estará en la clavija “hembra”, nunca en el “macho”, para evitar contactos directos.

Los interruptores diferenciales se instalarán de acuerdo con las siguientes sensibilidades:
\begin{itemize}[label=--]
    \item\SI{300}{\milli\ampere}. Alimentación a la maquinaria.
    \item\SI{30}{\milli\ampere}. Alimentación a la maquinaria como mejora del nivel de seguridad.
    \item\SI{30}{\milli\ampere}. Para las instalaciones eclécticas de alumbrado.
\end{itemize}
Las partes metálicas de todo equipo ecléctico dispondrán de toma de tierra.

El neutro de la instalación estará puesto a tierra.

La toma de tierra se efectuará a través de la pica o placa de cada cuadro general.

El hilo de toma de tierra, siempre estará protegido con macarrón en colores amarillo y verde. Se prohíbe expresamente utilizarlo para otros usos.

La iluminación mediante portátiles cumplirá la siguiente norma:
\begin{itemize}[label=--]
    \item Portalámparas estanco de seguridad con manto aislante, rejilla protectora de la bombilla dotada de gancho de cuelgue a la pared, manguera antihumedad, clavija de conexión normalizada.
    \item La iluminación de los tajos se situará a una altura en torno a los 2 m. medidos desde la superficie de apoyo de los operarios en el puesto de trabajo.
    \item Las zonas de paso de la obra, estarán permanentemente iluminadas evitando rincones oscuros.  
\end{itemize}
No se permitirá las conexiones a tierra a través de conductores de agua.
	
No se permitirá el tránsito de carretillas y personas sobre mangueras eléctricas.

No se permitirá el tránsito bajo líneas eléctricas con elementos longitudinales transportados a hombros (pértigas, reglas, escaleras de mano...). La inclinación de la pieza puede llegar a producir contacto eléctrico.
\section{Señalización}
\label{sec:sys-señalizacion}
Se realizará la señalización oportuna según el tipo de trabajo que se esté realizando, la fase de ejecución y el lugar del mismo. Las señalizaciones serán temporales, durarán el tiempo que se prolongue los trabajos. Serán de tipo: triángulos con hombres trabajando, cintas, banderolas...

Cuando por cruzamientos sea necesario advertir de los límites de velocidad y altura, estrechamiento de la calzada, etc. se colocarán estas señales antes y después del lugar de trabajo, a la distancia reglamentaria para cada tipo de carretera.

La señalización fija que debe llevar las instalaciones eléctricas estarán prescritas en el Reglamento para Líneas Eléctricas de Alta Tensión. Dicha señalización previene del riesgo que supone la electricidad , prohibiendo tocar los conductores y apoyos. Esta señalización se coloca en los apoyos.
